% Generated draft: requires mathematical and source review.
\documentclass[11pt]{article}
\usepackage[margin=25mm]{geometry}
\usepackage{fontspec}
\setmainfont{Latin Modern Roman}
\usepackage{amsmath,amssymb,mathtools,unicode-math}
\setmathfont{Latin Modern Math}
\usepackage{enumitem,needspace,xurl,hyperref}
\setlength{\parindent}{0pt}
\setlength{\parskip}{5pt}
\setlength{\emergencystretch}{3em}
\newcommand{\N}{\mathbb{N}}
\newcommand{\Z}{\mathbb{Z}}
\newcommand{\Q}{\mathbb{Q}}
\newcommand{\R}{\mathbb{R}}
\newcommand{\C}{\mathbb{C}}
\newcommand{\F}{\mathbb{F}}
\newcommand{\A}{\mathbb{A}}
\newcommand{\PP}{\mathbb{P}}
\newcommand{\E}{\mathbb{E}}
\newcommand{\eps}{\varepsilon}
\newcommand{\dd}{\,\mathrm{d}}
\begin{document}
\section*{UnsolvedMath ID: 3100058}
\textbf{Source problem code:} catalogue entry
\subsection{Definitions and mathematical statement}
Let \(S\subseteq\N_{>0}\). The source calls \(S\) sum-free when
\[
\forall s\in S\ \forall F\subseteq S\setminus\{s\},\qquad
0<|F|<\infty\ \Longrightarrow\ s\ne\sum_{a\in F}a.
\]
Thus summands are distinct elements of a finite subset; this is stronger than the usual condition forbidding only \(x+y=z\). Define the nonnegative reciprocal sum and extremal value by
\[
R(S)=\sup_{\substack{F\subseteq S\\ |F|<\infty}}\sum_{a\in F}\frac1a,
\qquad
R=\sup\{R(S):S\subseteq\N_{>0}\text{ satisfies the displayed condition}\}.
\]
The question is to determine \(R\). The source records \(2.064<R<4\), attributing the upper bound to Abbott and the lower bound to Levine--O'Sullivan. It asks for the supremum; attainment by an infinite or finite set is not part of its stated conclusion.

The supplied research summary labels the exact value open with partial progress, but its further numerical and attribution claims are tentative and conflict with the source definition. In particular, the positive odd integers fail the displayed condition since \(9=1+3+5\), and their reciprocal sum diverges; they cannot yield the finite value suggested in that summary. Its references to Cilleruelo--Hamma, Schoen, and a possible Schoen--Tomon upper bound near \(3.96\) lack precise citations and are retained only as unverified research notes, not added bounds or an identified maximizer.
\subsection{Short English statement}
Determine the largest possible reciprocal sum, in the sense of a supremum, for sets of positive integers in which no element is a sum of distinct other elements. The source gives \(2.064<R<4\); the supplied later-status summary needs correction and verification.
\subsection{Sources}
\begin{itemize}
\item Local UnsolvedMath registry, permanent ID \(3100058\), exact source code \texttt{catalogue entry}; problem attributed to Erdős.
\item Cooper, \emph{Combinatorial Problems I Like}, October 2020 snapshot, Combinatorial Number Theory, source-order bullet 58; \url{https://people.math.sc.edu/cooper/combprob.html}. Recorded access 2026-07-29; HTML extraction.
\item Source statement attributes \(R<4\) to Abbott and \(R>2.064\) to Levine--O'Sullivan. Complete bibliographic details are not supplied.
\item Supplied research summary mentions Cilleruelo--Hamma, Schoen, an ``Abbott--\ldots'' improvement, and ``Schoen--Tomon?''; these tentative attributions and numerical claims have no accompanying precise references.
\end{itemize}
\end{document}
